\documentclass[11pt]{article}

\usepackage[margin=1in]{geometry}
\usepackage{amsmath,amssymb,amsthm}
\usepackage{booktabs}
\usepackage{natbib}
\usepackage{url}

\newtheorem{definition}{Definition}
\newtheorem{proposition}{Proposition}
\newtheorem{remark}{Remark}

\newcommand{\R}{\mathbb{R}}
\newcommand{\argmin}{\operatorname*{arg\,min}}
\newcommand{\diag}{\operatorname{diag}}
\newcommand{\vech}{\operatorname{vech}}
\newcommand{\vecl}{\operatorname{vecl}}
\newcommand{\trans}{^{\mathsf T}}
\newcommand{\magmaan}{\texttt{magmaan}}
\newcommand{\lavaan}{\texttt{lavaan}}
\newcommand{\SNLLS}{\mathrm{SNLLS}}
\newcommand{\ULS}{\mathrm{ULS}}
\newcommand{\DWLS}{\mathrm{DWLS}}
\newcommand{\WLS}{\mathrm{WLS}}

\title{Separable Least Squares for Ordinal Structural Equation Models}
\author{}
\date{\today}

\begin{document}
\maketitle

\begin{abstract}
Separable nonlinear least squares (SNLLS) profiles the conditionally linear
parameters in fixed-weight least-squares estimation of continuous structural
equation models. Ordinal SEM is not a plug-in application of the same
reduction. Its least-squares target contains thresholds and latent-response
correlations, and the full weight matrix couples the two blocks. This note
sketches an ordinal SNLLS formulation. Thresholds are represented as an
ordinal measurement layer outside the LISREL covariance block, free thresholds
can be profiled analytically, and full WLS threshold profiling changes the
effective correlation weight through a Schur complement. The resulting
correlation problem can reuse the spirit of continuous SNLLS, but only after
the ordinal moment vector, threshold constraints, and response-scale
parameterization are made explicit.
\end{abstract}

\section{Motivation}
\label{sec:motivation}

Fixed-weight least-squares estimation of continuous structural equation models
compares sample means and covariances with the corresponding model-implied
moments. In the usual LISREL notation \citep{joreskog1996lisrel},
\begin{align}
  \Sigma(\theta)
    &=
    \Lambda(I-B)^{-1}\Psi(I-B)^{-\mathsf T}\Lambda\trans+\Theta, \\
  \mu(\theta)
    &=
    \nu+\Lambda(I-B)^{-1}\alpha.
\end{align}
With sample moments \(s=(\bar x\trans,\vech(S)\trans)\trans\) and model moments
\(m(\theta)=(\mu(\theta)\trans,\vech\{\Sigma(\theta)\}\trans)\trans\), the
criterion is
\begin{equation}
  F(\theta)
  =
  \frac{1}{2}\{s-m(\theta)\}\trans W\{s-m(\theta)\}.
  \label{eq:continuous-ls}
\end{equation}
SNLLS applies the variable-projection idea of
\citet{golub1973differentiation,golub2003separable}: for fixed directed
parameters, the remaining covariance and mean parameters enter
\(m(\theta)\) affinely, so they can be recovered by an inner linear
least-squares solve. This gives a smaller outer optimization over the directed
block only \citep{kreiberg2021cfa,kreiberg2023sem,ernst2023trek}.

Ordinal SEM changes the object being fitted. The observed categorical
variables are treated as coarsened versions of latent continuous responses
\citep{muthen1984categorical}. For an ordinal variable \(Y_j\) with
\(K_j\) categories, introduce a latent response \(Y_j^\star\) and thresholds
\[
  -\infty=\tau_{j0}<\tau_{j1}<\cdots<\tau_{j,K_j-1}<\tau_{j,K_j}=\infty,
  \qquad
  Y_j=k
  \Longleftrightarrow
  \tau_{j,k-1}<Y_j^\star\leq\tau_{jk}.
\]
The polychoric stage estimates thresholds and latent-response correlations
\citep{olsson1979polychoric}. Thus the observed least-squares moment vector is
not merely \(\vech(S)\). It is an augmented vector
\begin{equation}
  s_o =
  \begin{bmatrix}
    \hat\tau \\
    \hat\rho
  \end{bmatrix},
  \label{eq:ordinal-sample-moments}
\end{equation}
where \(\hat\tau\) stacks all finite thresholds and \(\hat\rho\) stacks the
off-diagonal latent-response correlations. This is why ordinal SNLLS is not a
simple call to the continuous SNLLS routine with \(S\) replaced by a
polychoric matrix.

The purpose of this note is to state the difference cleanly. The thresholds
belong to the measurement model but not to \(\Sigma(\theta)\). They can be
represented as lavaan-style partable rows, for example with the threshold
operator \texttt{|}, while the LISREL block continues to define the
latent-response covariance or correlation structure. The ordinal least-squares
criterion binds these pieces together through the moment vector and its weight.

\section{Ordinal moment criterion}
\label{sec:criterion}

Let \(R^\star(\theta)\) be the model-implied latent-response correlation
matrix for the ordinal observed variables. In a delta-style parameterization
the diagonal of \(R^\star(\theta)\) is fixed to one; in a covariance
parameterization, \(R^\star(\theta)=D(\theta)^{-1/2}
\Sigma^\star(\theta)D(\theta)^{-1/2}\) with
\(D(\theta)=\diag\{\Sigma^\star(\theta)\}\). Write
\[
  \rho(\theta)=\vecl\{R^\star(\theta)\},
\]
where \(\vecl\) stacks the strict lower triangle. The model-implied ordinal
moment vector is
\begin{equation}
  m_o(\theta,\gamma)
  =
  \begin{bmatrix}
    \tau(\gamma) \\
    \rho(\theta)
  \end{bmatrix},
  \label{eq:ordinal-model-moments}
\end{equation}
where \(\gamma\) collects threshold parameters. The fixed-weight ordinal
least-squares criterion is
\begin{equation}
  F_o(\theta,\gamma)
  =
  \frac{1}{2}
  \{s_o-m_o(\theta,\gamma)\}\trans
  W_o
  \{s_o-m_o(\theta,\gamma)\}.
  \label{eq:ordinal-ls}
\end{equation}
The choices \(\ULS\), \(\DWLS\), and \(\WLS\) differ only through \(W_o\).
For \(\ULS\), \(W_o=I\). For \(\DWLS\), \(W_o\) is diagonal. For full
\(\WLS\), \(W_o\) is the inverse, or generalized inverse, of the asymptotic
covariance matrix of \((\hat\tau\trans,\hat\rho\trans)\trans\)
\citep{browne1984adf,bollen1989sem}. In practice, categorical SEM software
often estimates parameters with a diagonal weight while retaining the full
weight machinery for robust standard errors and test statistics
\citep{flora2004ordinal,li2016ordinal,rosseel2012lavaan}.

\begin{remark}[Why polychorics are not enough]
The polychoric correlation matrix gives \(\hat\rho\), and the asymptotic
moment covariance gives the raw material for \(W_o\). But the continuous
criterion \eqref{eq:continuous-ls} is built on covariance cells, including
diagonals, while the ordinal criterion \eqref{eq:ordinal-ls} is built on
thresholds and off-diagonal correlations. The model side must therefore expose
both \(\tau(\gamma)\) and \(\rho(\theta)\), and the weight matrix must be
conformed to this augmented vector.
\end{remark}

\section{Threshold profiling}
\label{sec:threshold-profiling}

Thresholds are the part of the ordinal problem that is most visibly absent
from the continuous LISREL covariance block. They are also the easiest part to
profile. Suppose the threshold model is affine:
\begin{equation}
  \tau(\gamma)=\tau_0+H\gamma.
  \label{eq:threshold-affine}
\end{equation}
The matrix \(H\) encodes free thresholds, equality labels, fixed thresholds,
and invariance restrictions. With
\[
  d_\tau=\hat\tau-\tau_0,
  \qquad
  d_\rho(\theta)=\hat\rho-\rho(\theta),
  \qquad
  W_o=
  \begin{bmatrix}
    W_{\tau\tau} & W_{\tau\rho}\\
    W_{\rho\tau} & W_{\rho\rho}
  \end{bmatrix},
\]
the criterion can be written as a quadratic function of \(\gamma\):
\[
  F_o(\theta,\gamma)
  =
  \frac{1}{2}
  \begin{bmatrix}
    d_\tau-H\gamma \\
    d_\rho(\theta)
  \end{bmatrix}\trans
  W_o
  \begin{bmatrix}
    d_\tau-H\gamma \\
    d_\rho(\theta)
  \end{bmatrix}.
\]
For fixed \(\theta\), the profiled threshold estimate is
\begin{equation}
  \hat\gamma(\theta)
  =
  \{H\trans W_{\tau\tau}H\}^{+}
  H\trans
  \{W_{\tau\tau}d_\tau+W_{\tau\rho}d_\rho(\theta)\},
  \label{eq:threshold-profile-general}
\end{equation}
where \(+\) denotes the Moore--Penrose inverse. If \(H\) has full column rank
and \(W_{\tau\tau}\) is positive definite on its span, this is an ordinary
weighted least-squares solve \citep{golub2013matrix,magnus1999matrix}.

\begin{proposition}[Free-threshold profiling]
\label{prop:free-threshold}
If every finite threshold is freely estimated, so that \(H=I\) and
\(\tau_0=0\), profiling thresholds out of \eqref{eq:ordinal-ls} gives the
correlation-only criterion
\begin{equation}
  F_{o,p}(\theta)
  =
  \frac{1}{2}d_\rho(\theta)\trans
  \left(
    W_{\rho\rho}
    -
    W_{\rho\tau}W_{\tau\tau}^{-1}W_{\tau\rho}
  \right)
  d_\rho(\theta).
  \label{eq:threshold-schur}
\end{equation}
\end{proposition}

\noindent
The result is immediate from minimizing the quadratic over the unrestricted
threshold residual. The profiled threshold residual is not generally zero:
\[
  \hat\tau-\hat\gamma(\theta)
  =
  -W_{\tau\tau}^{-1}W_{\tau\rho}d_\rho(\theta).
\]
It is zero only when the threshold and correlation blocks are uncoupled by the
weight.

\begin{remark}[Full WLS versus DWLS and ULS]
For \(\ULS\) and \(\DWLS\), \(W_{\tau\rho}=0\), so free thresholds profile to
their sample estimates and the remaining fit is a correlation-only least
squares problem with the identity or diagonal correlation weights. For full
\(\WLS\), the threshold-correlation block is generally nonzero, so profiling
thresholds changes the correlation weight by the Schur complement in
\eqref{eq:threshold-schur}. If \(W_o=\Gamma_o^{-1}\), where
\(\Gamma_o\) is partitioned conformably with \(s_o\), the Schur complement in
\eqref{eq:threshold-schur} equals the inverse of the marginal covariance block
for \(\hat\rho\). Thus ``use the polychoric matrix and the lower-right block of
the inverse weight'' is not the same operation as profiling free thresholds.
\end{remark}

Constrained thresholds are the case that makes the ordinal layer visible. A
fixed threshold, an equality label across items, or a multi-group threshold
invariance restriction reduces the span of \(H\). Then threshold residuals
cannot be eliminated, and the profiled objective keeps the part of
\(d_\tau\) orthogonal to the threshold design, as well as its coupling with
\(d_\rho(\theta)\). This is why threshold syntax is not merely display
metadata: it defines a linear model for the threshold block.

\section{Ordinal SNLLS after threshold profiling}
\label{sec:ordinal-snlls}

After thresholds have either been profiled or retained explicitly, the
remaining nonlinear problem concerns the model-implied latent-response
correlations. The continuous SNLLS split still suggests a useful reduction,
but it must be stated on the ordinal moment scale.

Let \(b\) collect the directed parameters, such as free entries of
\(\Lambda\) and \(B\), and let \(c\) collect the variance, covariance, and mean
parameters that are conditionally linear in the continuous LISREL moment map.
For fixed \(b\), the latent-response covariance matrix has the affine form
\begin{equation}
  \sigma^\star(b,c)=\sigma^\star_0(b)+M^\star(b)c,
  \label{eq:latent-response-affine}
\end{equation}
where \(\sigma^\star\) is a vectorized covariance representation. Continuous
SNLLS applies directly to this form. Ordinal SNLLS applies directly only if
the ordinal parameterization preserves the relevant affine map from \(c\) to
the fitted ordinal moments.

The clean case is a delta-style all-ordinal model in which the
latent-response variances are fixed to one and the fitted moments are the
strict lower-triangular covariance entries, which are also correlations. Then
\[
  \rho(b,c)=\rho_0(b)+M_\rho(b)c,
\]
and, after free-threshold profiling, the ordinal criterion is again separable:
\begin{equation}
  F_{\SNLLS,o}(b)
  =
  \min_c
  \frac{1}{2}
  \{d_\rho-\rho_0(b)-M_\rho(b)c\}\trans
  W_{\rho\mid\tau}
  \{d_\rho-\rho_0(b)-M_\rho(b)c\},
  \label{eq:ordinal-snlls-clean}
\end{equation}
where \(W_{\rho\mid\tau}\) is the identity, a diagonal matrix, or the
threshold-profiled WLS weight from \eqref{eq:threshold-schur}. Conditional on
\(b\), the inner solution is again a linear least-squares solve. This is the
closest ordinal analogue of continuous SNLLS.

The less clean case is an unconstrained covariance parameterization in which
the model-implied correlations are obtained by standardizing a covariance
matrix:
\[
  \rho_{ij}(b,c)
  =
  \frac{\Sigma^\star_{ij}(b,c)}
       {\{\Sigma^\star_{ii}(b,c)\Sigma^\star_{jj}(b,c)\}^{1/2}}.
\]
Even if \(\Sigma^\star(b,c)\) is affine in \(c\), the standardized correlation
map is not. Then the continuous SNLLS profile does not carry over without an
additional scale parameterization or a constrained inner solve that enforces
the response variances. The same issue appears in mixed continuous-ordinal
models, where some observed moments are covariances and others are
latent-response correlations.

\begin{proposition}[Sufficient condition for ordinal separability]
\label{prop:ordinal-separable}
Suppose (i) thresholds follow the affine model \eqref{eq:threshold-affine},
(ii) threshold profiling is performed by \eqref{eq:threshold-profile-general}
or thresholds are included as part of the inner linear block, and
(iii) conditional on directed parameters \(b\), the fitted non-threshold
ordinal moments are affine in a block \(c\). Then the fixed-weight ordinal
least-squares problem is separable in \((b,c,\gamma)\). The outer objective is
obtained by one linear least-squares solve for thresholds and one linear
least-squares solve for \(c\) at each trial value of \(b\), or by a single
combined linear solve when the two blocks are retained together.
\end{proposition}

\begin{remark}
Condition (iii) is the substantive ordinal restriction. It is automatic for
continuous covariance moments and for ordinal correlation moments under a
parameterization that fixes the response variances before the least-squares
map is formed. It is not automatic for a generic covariance-to-correlation
standardization.
\end{remark}

\section{Estimator variants}
\label{sec:variants}

The same notation separates four estimators that are easily conflated.

\begin{description}
\item[ULS.]
Use \(W_o=I\). With free thresholds, the threshold block profiles away and
the outer problem fits only \(\hat\rho-\rho(\theta)\) by unweighted least
squares. This is the lightest ordinal SNLLS target.

\item[DWLS.]
Use a diagonal \(W_o\). With free thresholds, the threshold block again
profiles away, but the correlation residuals receive diagonal weights from
the asymptotic variances of the estimated polychoric correlations. This is the
natural analogue of the parameter-estimation step in common categorical SEM
workflows.

\item[Profiled WLS.]
Start from the full \(W_o\), profile free thresholds, and fit correlations
with \(W_{\rho\mid\tau}\) from \eqref{eq:threshold-schur}. This is cheaper
than carrying the full threshold-correlation residual through the outer
optimizer, but it is not the same as DWLS and not the same as using the
lower-right block of \(W_o\).

\item[Full ordinal WLS.]
Retain thresholds and correlations in one residual vector and use the full
weight \(W_o\). This is the direct minimum-distance estimator for the
augmented ordinal moment vector. It is the most faithful to the moment
criterion, but it is also where weight construction and inversion dominate
the computational story.
\end{description}

This classification gives a natural benchmarking design for a later paper
revision: separate construction of thresholds, polychorics, \(\Gamma_o\), and
the inverted weights from the repeated optimizer cost, then compare full
optimization with ordinal SNLLS under each weight choice. No such simulations
are part of the present note.

\section{Implementation implications}
\label{sec:implementation}

The representation should mirror the statistical decomposition. The LISREL
matrix representation continues to own \(\Sigma^\star(\theta)\) and
\(\mu(\theta)\). Thresholds live beside it as ordinal measurement rows, not
inside the covariance formula. The sample side owns the estimated thresholds,
polychoric correlations, and their asymptotic covariance. The fitting adapter
constructs the ordinal residual vector and applies the selected weight.

For \magmaan{}, this suggests an implementation boundary rather than a source
change in this note:
\[
  \text{partable threshold rows}
  \quad+\quad
  \text{ordinal sample statistics}
  \quad\longrightarrow\quad
  \text{ordinal SNLLS adapter}.
\]
The adapter would expose three explicit choices:
\begin{enumerate}
\item whether thresholds are free, fixed, or constrained through a design
  matrix \(H\);
\item whether the non-threshold moment map is already affine in the profiled
  block or requires response-scale handling;
\item whether the weight is \(\ULS\), \(\DWLS\), profiled \(\WLS\), or full
  \(\WLS\).
\end{enumerate}
This is deliberately not a change to the core LISREL evaluator. The point is
that ordinal SEM adds a measurement side layer and a different moment vector,
not a new formula for \(\Sigma(\theta)\).

\section{Open problems}
\label{sec:open}

Several details decide whether ordinal SNLLS is merely neat or actually useful.
First, the response-scale parameterization must be chosen so that the inner
linear solve is well-defined for ordinary CFA and SEM examples. Second,
threshold constraints need a lavaan-compatible interpretation: right-hand
threshold slots such as \texttt{t1}, \texttt{t2}, and \texttt{t3} should be
validated as ordered slots, while labels before \texttt{*} define equality
constraints. Third, multi-group threshold invariance is just a structured
choice of \(H\), but it needs to be materialized before the profiler sees the
problem. Fourth, the full-WLS and profiled-WLS distinction should be tested
numerically, because the Schur-complement weight may behave differently from
both DWLS and the unprofiled full criterion in small samples.

\section{Discussion}
\label{sec:discussion}

Ordinal SNLLS is a separate methodological object from continuous SNLLS. The
continuous reduction profiles parameters in a covariance and mean moment map.
The ordinal reduction must first decide what to do with thresholds and with
the covariance-to-correlation scale of the latent responses. Once those
decisions are explicit, the attractive case is simple: free thresholds profile
analytically, ULS and DWLS collapse to correlation-only criteria, and full WLS
collapses through a Schur-complement weight. The remaining correlation fit can
then use the same variable-projection machinery as continuous SNLLS when the
response-scale parameterization preserves conditional linearity.

That is the small but real contribution. Polychoric correlations and Gamma are
the correct ingredients, but the recipe is different: the estimator lives on an
augmented threshold-correlation moment vector. Treating the threshold layer
explicitly is what makes ordinal SNLLS a method rather than an implementation
trick.

\bibliographystyle{plainnat}
\bibliography{ordinal-snlls}

\end{document}
